\chapter[Lezione 1 --- 29 settembre 2026]{Lezione 1: numeri naturali, ordinamento e principio di induzione}
\label{chap:lezione-01}

\section{Numeri naturali e le loro operazioni}
\label{sec:lez01-numeri-naturali}

Si considera l'insieme dei numeri naturali
\[
  \N=\{0,1,2,3,\ldots\},
\]
adottando quindi la convenzione \(0\in\N\).

Su \(\N\) sono definite le operazioni di somma e prodotto. Entrambe sono
operazioni binarie, ossia applicazioni che a ogni coppia ordinata
\((a,b)\in\N\times\N\) associano un elemento di \(\N\):
\[
  +\colon\N\times\N\to\N,
  \qquad
  (a,b)\mapsto a+b,
\]
\[
  \cdot\colon\N\times\N\to\N,
  \qquad
  (a,b)\mapsto ab.
\]
L'insieme \(\N\times\N\) costituisce pertanto l'insieme di partenza,
mentre \(\N\) è l'insieme di arrivo delle due operazioni.

Le proprietà fondamentali della somma e del prodotto introdotte di seguito
vengono assunte come \emph{assiomi} e costituiscono il punto di
partenza per i successivi sviluppi, senza essere dimostrate.

Per ogni \(a,b,c\in\N\), la somma è commutativa, ossia l'ordine degli
addendi non modifica il risultato, e associativa, ossia il modo in cui
gli addendi vengono raggruppati non modifica la somma:
\[
  a+b=b+a,
  \qquad
  (a+b)+c=a+(b+c).
\]
La somma ammette un unico elemento neutro \(e_0\), caratterizzato
dalla proprietà
\[
  \exists!\,e_0\in\N
  \quad\text{tale che}\quad
  \forall a\in\N,\qquad a+e_0=a.
\]
Tale elemento è
\[
  e_0=0.
\]

Per ogni \(a,b,c\in\N\), il prodotto è commutativo, ossia l'ordine dei
fattori non modifica il risultato, e associativo, ossia il modo in cui i
fattori vengono raggruppati non modifica il prodotto:
\[
  ab=ba,
  \qquad
  (ab)c=a(bc).
\]
Il prodotto ammette un unico elemento neutro \(e_1\), caratterizzato
dalla proprietà
\[
  \exists!\,e_1\in\N
  \quad\text{tale che}\quad
  \forall a\in\N,\qquad ae_1=a.
\]
Tale elemento è
\[
  e_1=1.
\]

\clearpage

\section{Ordinamento totale dei numeri naturali}
\label{sec:lez01-ordine-naturali}

Per \(a,b\in\N\), si definisce la relazione \(\le\) nel modo seguente.

\begin{definition}[Ordine naturale]
\label{def:ordine-naturali}
\[
  a\le b
  \quad\Longleftrightarrow\quad
  \exists c\in\N:\ a+c=b.
\]
\end{definition}

La relazione \(a\le b\) esprime quindi il fatto che è possibile ottenere
\(b\) aggiungendo ad \(a\) un opportuno numero naturale. Tale relazione
soddisfa quattro proprietà fondamentali.

\textbf{Riflessività.}
Ogni numero naturale è minore o uguale a se stesso; per ogni \(a\in\N\)
vale
\[
  a\le a.
\]
Infatti, nella \cref{def:ordine-naturali} è sufficiente scegliere \(c=0\),
poiché
\[
  a+0=a.
\]

\textbf{Antisimmetria.}
Se due numeri naturali sono ciascuno minore o uguale dell'altro, allora
coincidono. Per ogni \(a,b\in\N\),
\[
  a\le b
  \quad\text{e}\quad
  b\le a
  \quad\Longrightarrow\quad
  a=b.
\]

\textbf{Transitività.}
Se un numero naturale è minore o uguale a un secondo numero e quest'ultimo
è minore o uguale a un terzo, allora il primo è minore o uguale al terzo.
Per ogni \(a,b,c\in\N\),
\[
  a\le b
  \quad\text{e}\quad
  b\le c
  \quad\Longrightarrow\quad
  a\le c.
\]

\textbf{Dicotomia.}
Ogni coppia di numeri naturali è confrontabile: per ogni \(a,b\in\N\)
vale
\[
  a\le b
  \quad\text{oppure}\quad
  b\le a.
\]
Le due possibilità non sono necessariamente esclusive: quando \(a=b\),
entrambe le relazioni risultano vere.

Le proprietà precedenti caratterizzano la relazione \(\le\) come un
ordinamento totale su \(\N\).

La relazione d'ordine è inoltre compatibile con le operazioni di somma e
prodotto.

\begin{proposition}[Compatibilità dell'ordine]
\label{prop:compatibilita-ordine}
Per ogni \(a,b,c\in\N\), se \(a\le b\), allora
\[
  a+c\le b+c.
\]
Inoltre, se \(c>0\), vale
\[
  ac\le bc.
\]
\end{proposition}

La prima proprietà afferma che l'aggiunta dello stesso numero naturale ai
due membri preserva la relazione d'ordine; la seconda esprime l'analoga
compatibilità rispetto alla moltiplicazione per un numero naturale
strettamente positivo.

Se \(c=0\), si ha \(ac=bc=0\), per cui la disuguaglianza non stretta
rimane comunque valida.

\clearpage

\section{Richiami su insiemi e dimostrazioni}
\label{sec:lez01-insiemi-dimostrazioni}

Siano \(A\) e \(B\) due insiemi. Tra le operazioni fondamentali sugli
insiemi si considerano l'intersezione, che raccoglie gli elementi comuni,
e l'unione, che raccoglie gli elementi appartenenti ad almeno uno dei due:

\[
  A\cap B=\{x:x\in A\text{ e }x\in B\},\qquad
  A\cup B=\{x:x\in A\text{ oppure }x\in B\}.
\]

La \cref{fig:venn-operazioni} rappresenta le regioni corrispondenti:

\begin{figure}[htbp]
  \centering
  \begin{tikzpicture}[unipd/diagram]
  \begin{scope}
    \clip (0,0) circle (1);
    \fill[unipd-accent!25] (1.2,0) circle (1);
  \end{scope}
  \draw[unipd/edge] (0,0) circle (1);
  \draw[unipd/edge] (1.2,0) circle (1);
  \node at (-0.45,0) {$A$};
  \node at (1.65,0) {$B$};
  \node at (0.6,-1.35) {$A\cap B$};
  \begin{scope}[xshift=5cm]
    \fill[unipd-accent!25] (0,0) circle (1);
    \fill[unipd-accent!25] (1.2,0) circle (1);
    \draw[unipd/edge] (0,0) circle (1);
    \draw[unipd/edge] (1.2,0) circle (1);
    \node at (-0.45,0) {$A$};
    \node at (1.65,0) {$B$};
    \node at (0.6,-1.35) {$A\cup B$};
  \end{scope}
\end{tikzpicture}

  \caption{Intersezione (a sinistra) e unione (a destra) di due insiemi.}
  \label{fig:venn-operazioni}
\end{figure}

Accanto alla notazione insiemistica, è necessario fissare il linguaggio
impiegato negli enunciati matematici e nelle relative dimostrazioni.

Un teorema ha la forma
\[
  \text{IPOTESI}\Longrightarrow\text{TESI}.
\]
Una dimostrazione deduce la tesi dalle ipotesi mediante passaggi logicamente
validi. Per proposizioni \(A\) e \(B\), l'implicazione
\(A\Longrightarrow B\) asserisce che, se \(A\) è vera, allora è vera
anche \(B\). Nella logica classica essa è equivalente alla propria
contronominale:
\[
  A\Rightarrow B\quad\Longleftrightarrow\quad\neg B\Rightarrow\neg A.
\]
La dimostrazione per assurdo è distinta: mantenendo valide le ipotesi, si
assume la negazione della tesi e da tali assunzioni si ricava una
contraddizione, ossia una proposizione \(C\) e, contemporaneamente, la sua
negazione \(\neg C\).

\section{Principio di induzione}
\label{sec:lez01-induzione}

Il principio di induzione consente di dimostrare che una determinata
proprietà vale per tutti i numeri naturali a partire da un valore iniziale
\(p\). Si verifica anzitutto la proprietà per \(p\) e si dimostra che la
sua validità per un generico naturale \(n\) implica la validità per il
successivo \(n+1\). Ne segue così la propagazione della proprietà da
\(p\) a \(p+1\), quindi a \(p+2\), e così via.

\begin{theorem}[Principio di induzione]
\label{thm:principio-induzione}
Siano \(A\subseteq\N\) e \(p\in\N\). Se
\[
  p\in A
  \qquad\text{e}\qquad
  \forall n\in\N,\quad n\in A\Longrightarrow n+1\in A,
\]
allora
\[
  \{n\in\N:n\ge p\}\subseteq A.
\]
\end{theorem}

\clearpage

La prima condizione è
\[
  p\in A.
\]
Essa costituisce il \emph{passo base} e fissa il punto iniziale
dell'induzione. La seconda condizione è
\[
  n\in A\Longrightarrow n+1\in A.
\]
Essa costituisce il \emph{passo induttivo} e permette di estendere
l'appartenenza ad \(A\) da un naturale al suo successivo.

Se \(P(n)\) è una proprietà dipendente da \(n\) e si intende dimostrarne
la validità per ogni \(n\ge p\), si può porre
\[
  A=\{n\in\N : n\ge p \text{ e } P(n)\text{ è vera}\}.
\]
Dimostrare che \(P(n)\) vale per ogni \(n\ge p\) equivale allora a
verificare
\[
  P(p)
\]
e a dimostrare, per un generico \(n\ge p\),
\[
  P(n)\Longrightarrow P(n+1).
\]

Vediamo ora due applicazioni del principio di induzione:
nell'\cref{ex:somma-primi-naturali} dimostreremo la formula per la somma
dei primi \(n\) naturali positivi, mentre nell'\cref{ex:somma-geometrica-finita}
considereremo la somma geometrica finita.

\begin{example}[Somma dei primi \(n\) naturali positivi]
\label{ex:somma-primi-naturali}
Per ogni \(n\in\N\) con \(n\ge1\) vale
\[
  \sum_{i=1}^{n}i=\frac{n(n+1)}2.
\]
La formula fornisce quindi un'espressione chiusa per la somma
\[
  1+2+\cdots+n
\]
dei primi \(n\) numeri naturali positivi.
\end{example}

\begin{proof}
Per applicare il principio di induzione si considera l'insieme dei naturali
per i quali la proprietà risulta vera:
\[
  A=\left\{n\in\N\,\middle|\,n\ge1,\;
    \sum_{i=1}^{n}i=\frac{n(n+1)}2\right\}.
\]
L'obiettivo è mostrare che ogni numero naturale maggiore o uguale a \(1\)
appartiene ad \(A\).

\emph{Passo base.}
Si verifica anzitutto la proprietà per \(n=1\). Si ha
\[
  \sum_{i=1}^{1}i=1
  \qquad\text{e}\qquad
  \frac{1(1+1)}2=1.
\]
I due membri coincidono; pertanto \(1\in A\) e il passo base è verificato.

\emph{Passo induttivo.}
Sia \(n\in A\). Per definizione di \(A\), si ha \(n\ge1\) e vale
l'ipotesi induttiva
\[
  \sum_{i=1}^{n}i=\frac{n(n+1)}2.
\]
Si deve dimostrare che anche il naturale successivo appartiene ad \(A\),
ossia che \(n+1\in A\). Poiché \(n+1\ge1\), la tesi induttiva consiste
nel verificare
\[
  \sum_{i=1}^{n+1}i=\frac{(n+1)(n+2)}2.
\]

Separando dalla sommatoria l'ultimo addendo e applicando l'ipotesi
induttiva, si ottiene
\begin{align*}
  \sum_{i=1}^{n+1}i
    &=\sum_{i=1}^{n}i+(n+1)\\
    &=\frac{n(n+1)}2+(n+1)\\
    &=\frac{n(n+1)+2(n+1)}2\\
    &=\frac{(n+1)(n+2)}2.
\end{align*}
L'ultima espressione coincide con quella richiesta dalla tesi induttiva;
pertanto \(n+1\in A\). Risulta quindi verificata l'implicazione
\[
  n\in A\Longrightarrow n+1\in A.
\]

Essendo verificati sia il passo base sia il passo induttivo, per il
\cref{thm:principio-induzione}, con valore iniziale \(p=1\), si conclude che
\[
  \{n\in\N:n\ge1\}\subseteq A.
\]
Di conseguenza,
\[
  \sum_{i=1}^{n}i=\frac{n(n+1)}2
\]
per ogni \(n\in\N\) con \(n\ge1\).
\end{proof}

\begin{example}[Somma geometrica finita]
\label{ex:somma-geometrica-finita}
Sia \(q\in\R\setminus\{1\}\). Per ogni \(n\in\N\) vale
\[
  \sum_{k=0}^{n}q^k=\frac{1-q^{n+1}}{1-q}.
\]
Con la convenzione \(q^0=1\), il membro sinistro rappresenta la somma
\[
  1+q+q^2+\cdots+q^n.
\]
La formula fornisce quindi un'espressione chiusa per la somma dei primi
\(n+1\) termini della progressione geometrica di ragione \(q\).

In particolare, per \(q=0\) tutti i termini con esponente positivo sono
nulli e si ha
\[
  \sum_{k=0}^{n}0^k
  =1+0+\cdots+0
  =1.
\]
D'altra parte,
\[
  \frac{1-0^{n+1}}{1-0}=1.
\]
Il caso \(q=0\) soddisfa quindi la formula; l'unico valore escluso è
\(q=1\), poiché annulla il denominatore.
\end{example}

\begin{proof}
Fissato \(q\in\R\setminus\{1\}\), si considera l'insieme dei naturali
per i quali la proprietà risulta vera:
\[
  A=\left\{n\in\N\,\middle|\,
    \sum_{k=0}^{n}q^k=\frac{1-q^{n+1}}{1-q}\right\}.
\]
L'obiettivo è mostrare che ogni numero naturale appartiene ad \(A\),
ossia che \(A=\N\).

\emph{Passo base.}
Si verifica anzitutto la proprietà per \(n=0\). Si ha
\[
  \sum_{k=0}^{0}q^k=q^0=1
\]
e, poiché \(q\neq1\),
\[
  \frac{1-q^{0+1}}{1-q}
  =\frac{1-q}{1-q}
  =1.
\]
I due membri coincidono; pertanto \(0\in A\) e il passo base è verificato.

\emph{Passo induttivo.}
Sia \(n\in A\). Per definizione di \(A\), vale l'ipotesi induttiva
\[
  \sum_{k=0}^{n}q^k=\frac{1-q^{n+1}}{1-q}.
\]
Si deve dimostrare che anche il naturale successivo appartiene ad \(A\),
ossia che \(n+1\in A\). La tesi induttiva consiste quindi nel verificare
\[
  \sum_{k=0}^{n+1}q^k=\frac{1-q^{n+2}}{1-q}.
\]

Separando dalla sommatoria l'ultimo addendo e applicando l'ipotesi
induttiva, si ottiene
\begin{align*}
  \sum_{k=0}^{n+1}q^k
    &=\sum_{k=0}^{n}q^k+q^{n+1}\\
    &=\frac{1-q^{n+1}}{1-q}+q^{n+1}\\
    &=\frac{1-q^{n+1}+q^{n+1}(1-q)}{1-q}\\
    &=\frac{1-q^{n+1}+q^{n+1}-q^{n+2}}{1-q}\\
    &=\frac{1-q^{n+2}}{1-q}.
\end{align*}
L'ultima espressione coincide con quella richiesta dalla tesi induttiva;
pertanto \(n+1\in A\). Risulta quindi verificata l'implicazione
\[
  n\in A\Longrightarrow n+1\in A.
\]

Essendo verificati sia il passo base sia il passo induttivo, per il
\cref{thm:principio-induzione}, con valore iniziale \(p=0\), si conclude che
\[
  A=\N.
\]
Poiché \(q\in\R\setminus\{1\}\) era arbitrario, si conclude che
\[
  \sum_{k=0}^{n}q^k=\frac{1-q^{n+1}}{1-q}
\]
per ogni \(n\in\N\) e per ogni \(q\in\R\setminus\{1\}\).
\end{proof}
